% !TeX root = ../../Book5.tex
% 纲要标题：### C.1 Lie 代数表示与 Lie 群表示的积分
% 纲要位置：工作记录/计划纲要.md，第 4873 行。
% #### C.1.1 微分表示与积分定理
% #### C.1.2 下降到商群的条件
% #### C.1.3 \texorpdfstring{\(SU(2)\) 与 \(SO(3)\) 的表示差别}{SU(2) 与 SO(3) 的表示差别}
% 节标题由计划纲要.md 中附录 C 的第 1 条要点整理。

\section{Lie 代数表示与 Lie 群表示的积分}
\label{b5:sec:C01}

% 纲要范围：
% * Lie 代数表示何时积分成 Lie 群表示

Lie 代数记录单位元附近的一阶变化。要把这种变化沿整个群延伸，还须检查群的连通性与绕闭路的影响。本节只讨论有限维表示；无限维空间上还要指定算子的定义域和拓扑，不能直接沿用这里的结论。

\subsection{微分表示与积分定理}
\label{b5:subsec:C01:01}

\begin{definition}\label{b5:def:lie-group-representation}
设 \(G\) 为实 Lie 群，\(V\) 为有限维复向量空间。光滑群同态 \(\pi:G\to\operatorname{GL}_{\mathbb C}(V)\) 称为 \(G\) 的有限维复表示。它在单位元处的微分
\(\mathrm{d}\pi:\operatorname{Lie}(G)\to\mathfrak{gl}_{\mathbb C}(V)\)
是实线性 Lie 代数同态。给定实线性 Lie 代数表示 \(\rho:\operatorname{Lie}(G)\to\mathfrak{gl}_{\mathbb C}(V)\)，若 \(\mathrm{d}\pi=\rho\)，就称 \(\pi\) 是 \(\rho\) 的\term[jifen-biaoshi]{积分表示}[integrated representation]。若 \(G\) 为复 Lie 群，则本节所说的复表示要求群同态全纯，其微分要求复线性。
\end{definition}

对一参数子群有
\[
\pi(\exp_G tx)=\exp\bigl(t\,\mathrm{d}\pi(x)\bigr).
\]
左边是一般线性群中的一参数子群，其单位元处导数为 \(\mathrm{d}\pi(x)\)；常微分方程 \(Y'=\mathrm{d}\pi(x)Y\)、\(Y(0)=I\) 的唯一解就是右边。微分保持括号可以从群交换子在二阶处的展开得到；同样的展开已经在矩阵 Lie 代数中出现。

\begin{theorem}[Lie 代数表示的积分]\label{b5:thm:lie-representation-integration}
设 \(G\) 为连通且单连通的实 Lie 群，\(\mathfrak g=\operatorname{Lie}(G)\)，\(V\) 为有限维复向量空间。每个实线性 Lie 代数表示 \(\rho:\mathfrak g\to\mathfrak{gl}_{\mathbb C}(V)\) 都唯一积分为光滑群表示 \(\pi:G\to\operatorname{GL}_{\mathbb C}(V)\)。微分给出有限维光滑复 \(G\)-表示与有限维实线性 \(\mathfrak g\)-表示的范畴等价。若 \(G\) 为连通且单连通的复 Lie 群，则每个有限维复线性 Lie 代数表示都唯一积分为全纯群表示，并给出相应的范畴等价。
\end{theorem}

\begin{proof}
设 \(\gamma:[0,1]\to G\) 是从单位元到 \(g\) 的分段光滑路径，令
\(a(t)=(\mathrm{d}L_{\gamma(t)^{-1}})\gamma'(t)\in\mathfrak g\)。沿这条路径解矩阵常微分方程
\[
 U'(t)=U(t)\rho\bigl(a(t)\bigr),\qquad U(0)=I.
\]
解在有限区间上存在唯一且可逆；例如行列式满足一个标量线性方程，初值为一。拟定义 \(\pi(g)=U(1)\)，关键是验证其与路径无关。

对保持端点的光滑同伦 \(\gamma(t,s)\)，再令
\(b(t,s)=(\mathrm{d}L_{\gamma(t,s)^{-1}})\partial_s\gamma(t,s)\)。左平移把各点不同的切空间都移到单位元处，故 \(a,b\) 可在同一个 Lie 代数中比较。矩阵群中，这就是 \(a=\gamma^{-1}\gamma_t\)、\(b=\gamma^{-1}\gamma_s\)，并且
\[
a_s=-ba+\gamma^{-1}\gamma_{ts},\qquad
b_t=-ab+\gamma^{-1}\gamma_{st}.
\]
两式相减即得下面的恒等式。一般 Lie 群中，它是左 Maurer--Cartan 形式\footnote{毛雷尔（Ludwig Maurer，1859-1927），德国数学家，研究连续变换群；嘉当已在本卷前文介绍。}的恒等式：
\[
 \partial_s a-\partial_t b=[a,b].
\]
若 \(U\) 真由路径端点决定，其沿同伦参数的变化应当满足 \(U_s=U\rho(b)\)。因此令 \(W=U_s-U\rho(b)\) 测量这一等式的误差。由于 \(\rho\) 保持括号，求导得
\[
\begin{aligned}
W_t&=U_s\rho(a)+U\rho(a_s)-U\rho(a)\rho(b)-U\rho(b_t)\\
&=W\rho(a)+U\bigl(\rho(b)\rho(a)-\rho(a)\rho(b)+\rho([a,b])\bigr)
=W\rho(a).
\end{aligned}
\]
端点固定使 \(b(0,s)=b(1,s)=0\)，故 \(W(0,s)=0\)，唯一性推出 \(W=0\)，进而 \(\partial_sU(1,s)=0\)。分段光滑同伦可逐段使用同一计算。Lie 群上的路径及同伦可以光滑化，而单连通性保证任意两条同端点路径同伦，所以定义良好。

把通向 \(g\) 的路径与左平移后的通向 \(h\) 的路径拼接，第二段的左对数导数仍为原路径的 \(a(t)\)。解方程便得
\(\pi(gh)=\pi(g)\pi(h)\)。在单位元附近取路径 \(t\mapsto\exp(tx)\)，则
\(\pi(\exp x)=e^{\rho(x)}\)，所以 \(\pi\) 光滑且 \(\mathrm{d}\pi=\rho\)。复 Lie 群情形中，这一局部公式是全纯的，再用群平移得到全群上的全纯性。

指数像含单位元邻域，而连通群由任何这样的邻域生成，故微分相同的两个表示必在全群上相同。对两个 Lie 代数表示 \(\rho_V,\rho_W\) 及其积分 \(\pi_V,\pi_W\)，若线性映射 \(T:V\to W\) 满足
\[
T\rho_V(x)=\rho_W(x)T\qquad(x\in\mathfrak g),
\]
则矩阵指数的幂级数给出
\[
Te^{\rho_V(x)}=e^{\rho_W(x)}T.
\]
由群的邻域生成性，得到 \(T\pi_V(g)=\pi_W(g)T\) 对所有 \(g\in G\) 成立。反向沿一参数子群对这一等式求微分，即得 Lie 代数作用的相容等式。于是微分既在对象上给出存在唯一性，又在同态上给出双射，得到范畴等价。上述路径构造使用常微分方程的存在唯一性及 Maurer--Cartan 恒等式；Lie 群同态版本见\cite[Theorem 3.41 and Theorem 4.3]{Kirillov2008LieGroups}。
\end{proof}

\subsection{下降到商群的条件}
\label{b5:subsec:C01:02}

连通 Lie 群 \(G\) 的普遍覆盖可以带上唯一使覆盖映射成为群同态的 Lie 群结构。记为
\(p:\widetilde G\to G\)，其核 \(D\) 是离散中心子群。离散性来自覆盖；中心性则因为对固定 \(d\in D\)，连续映射 \(g\mapsto gdg^{-1}\) 从连通的 \(\widetilde G\) 进入离散集合 \(D\)，只能恒为 \(d\)。

\begin{proposition}\label{b5:prop:representation-descent-cover}
设 \(G\) 为连通实 Lie 群，\(p:\widetilde G\to G\) 为作为 Lie 群同态的普遍覆盖，\(V\) 为有限维复向量空间，\(\rho:\operatorname{Lie}(G)\to\mathfrak{gl}_{\mathbb C}(V)\) 为实线性 Lie 代数表示。记 \(\widetilde\pi:\widetilde G\to\operatorname{GL}_{\mathbb C}(V)\) 为 \(\rho\circ\mathrm{d}p\) 的唯一积分。则 \(\rho\) 能积分为 \(G\) 的光滑复表示，当且仅当 \(\widetilde\pi(d)=1_V\) 对每个 \(d\in\ker p\) 成立。
\end{proposition}
\begin{proof}
若已有 \(G\) 表示，复合 \(p\) 后必零化覆盖核。反之，在覆盖核上平凡时，定义 \(\pi(g)=\widetilde\pi(\widetilde g)\)，其中 \(p(\widetilde g)=g\)。任意两个提升相差核元素，所以值不依赖提升。群乘法相容性由 \(\widetilde\pi\) 保证；覆盖局部是微分同胚，所以 \(\pi\) 光滑且具有原微分。
\end{proof}

这说明“Lie 代数表示可以积分”必须说明积分到哪个群。连通但非单连通的群可能丢掉一些表示；非连通群还需要补充连通分支群的作用，这些信息根本不在单位元切空间中。

\begin{example}\label{b5:exm:circle-integration}
对圆群 \(S^1=\mathbb R/2\pi\mathbb Z\)，一个 Lie 代数复表示由算子 \(A\) 给出，在普遍覆盖上的积分为 \(t\mapsto e^{tA}\)。下降条件为 \(e^{2\pi A}=I\)，等价于 \(A\) 可对角化且特征值属于 \(i\mathbb Z\)。确实，Jordan 块的半单部分必须满足 \(e^{2\pi\lambda}=1\)，其幂零部分 \(N\) 满足 \(e^{2\pi N}=I\)。但
\(e^{2\pi N}-I=N\cdot Q(N)\)，其中 \(Q(0)=2\pi\ne0\)，故 \(Q(N)\) 可逆，迫使 \(N=0\)。反向验证则直接代入。因此 \(A=i\alpha\) 的一维表示仅当 \(\alpha\in\mathbb Z\) 时下降；一个非零幂零 Jordan 块也不能下降。
\end{example}

\subsection{\texorpdfstring{\(SU(2)\) 与 \(SO(3)\) 的表示差别}{SU(2) 与 SO(3) 的表示差别}}
\label{b5:subsec:C01:03}

群 \(SU(2)\) 可写成
\[
\left\{\,
\begin{pmatrix}a&b\\-\overline b&\overline a\end{pmatrix}
\ \middle|\ |a|^2+|b|^2=1\,\right\}.
\]
作为流形它是三维球面，故单连通。其 Lie 代数的复化为 \(\mathfrak{sl}_2(\mathbb C)\)，因此前文的每个 \(L_m\) 唯一积分到 \(SU(2)\)。可以直接把它实现为自然二维表示的对称幂 \(\operatorname{Sym}^m\mathbb C^2\)。

\ProseParagraphBreak
\(SU(2)\) 在实三维空间 \(\mathfrak{su}_2\) 上的共轭作用保持正定内积 \(-\operatorname{tr}(XY)\)，给出 \(SU(2)\to SO(3)\)。它的微分是三维单 Lie 代数的伴随表示，单射且两边等维，故像含单位元邻域；\(SO(3)\) 连通使像为整个群。核由与全部 \(\mathfrak{su}_2\) 交换的矩阵组成，只能为 \(\{I,-I\}\)。因此它是二重覆盖。

\begin{proposition}\label{b5:prop:su2-so3-parity}
设 \(m\) 为非负整数，\(L_m\) 为最高权 \(m\) 的有限维简单左 \(\mathfrak{sl}_2(\mathbb C)\)-模。将它限制到实 Lie 代数 \(\mathfrak{su}_2\)，并积分为 \(SU(2)\) 的光滑复表示。对由伴随作用给出的二重覆盖 \(SU(2)\to SO(3)\)，其核为 \(\{I,-I\}\)，该积分表示能够下降为 \(SO(3)\) 表示，当且仅当 \(m\) 为偶数。
\end{proposition}
\begin{proof}
在 \(\operatorname{Sym}^m\mathbb C^2\) 上，\(-I\) 对每个次数 \(m\) 单项式按 \((-1)^m\) 作用。覆盖核由 \(-I\) 生成，所以\cref{b5:prop:representation-descent-cover}恰给出偶数条件。
\end{proof}

自然二维模是最小的反例：它是 \(SU(2)\) 的表示，却不是 \(SO(3)\) 的表示。三维伴随模对应 \(m=2\)，可以下降。同一个复化 Lie 代数因此并不意味着两个给定 Lie 群具有相同的表示。
